\documentclass[12pt]{article}
\usepackage{xeCJK}  

\usepackage{charter}
\usepackage{fullpage}
\usepackage[colorlinks=false,unicode]{hyperref}
\usepackage{ifthen}
\usepackage{comment}
\usepackage[title,titletoc]{appendix}
\usepackage{pagecolor}
\usepackage{amsmath}
\usepackage{amsfonts}
%\usepackage[normalem]{ulem}
\usepackage{siunitx}
\usepackage{amsthm}
\sisetup{per=slash, load=abbr}

\usepackage{pgfplots}
\usetikzlibrary{positioning}
\usetikzlibrary{fit}
\usetikzlibrary{snakes}
\usetikzlibrary{shapes.geometric}
\usetikzlibrary{patterns}
\usetikzlibrary{shapes,arrows,chains}
\usepgfplotslibrary{patchplots,colormaps}
\usetikzlibrary{calc}
\usetikzlibrary{positioning, fit}
\usetikzlibrary{backgrounds}
\usetikzlibrary{intersections}

\newcommand{\whitepaper}[1]{\begin{center}\fbox{\parbox{0.95\textwidth}{{\footnotesize
#1}}}\end{center}}

\newcommand{\pcolor}{orange!25}

\usepackage{setspace}
\usepackage{algorithm2e}
\bibliographystyle{ieeetr}

\usepackage{geometry}
\geometry{left=3cm,right=3cm,top=1.6cm,bottom=3cm,headheight=0pt,headsep=1.5em}
\usepackage{fancyhdr}
\pagestyle{fancy}
\chead{\includegraphics[scale=0.1]{../common/Nebulas.png}}  %在此处插入logo.pdf图片 图片靠左
\lhead{} % 页眉中间位置内容
\rhead{}
%\setlength{\topskip}{1em}

\usepackage{enumerate}
\setcounter{tocdepth}{4}
\setcounter{secnumdepth}{4}


\usepackage{indentfirst}


\setCJKmainfont[BoldFont = STSongti-SC-Bold]{STSongti-SC-Regular}
\setCJKfamilyfont{hei}{SIL-Hei-Med-Jian}		%宋体
\setCJKfamilyfont{song}{SimSun}		%宋体
\setCJKfamilyfont{kai}{Kaiti}		%楷体
\setCJKfamilyfont{fang}{song}	%仿宋
\setCJKfamilyfont{li}{song}			%隶书
\setCJKfamilyfont{you}{Yuanti}		%幼圆

\newcommand{\song}{\CJKfamily{song}}	%宋体
\newcommand{\hei}{\CJKfamily{hei}}	%黑体
\newcommand{\kai}{\CJKfamily{kai}}	%楷体
\newcommand{\fs}{\CJKfamily{fang}}	%仿宋
\newcommand{\li}{\CJKfamily{li}}		%隶书
\newcommand{\you}{\CJKfamily{you}}	%幼圆
\newcommand{\reffig}[1]{图\ref{#1}}
\newcommand{\refsec}[1]{\S \ref{#1}}

\onehalfspacing   % ----------设置1.5倍行距（可能有意义，待调整）

%\parindent=20pt  % -------------------首行缩进大小，英文分段就直接0pt了吧。
\setlength{\parindent}{2.1em}
\setlength{\parskip}{0.3\baselineskip}
\newcommand{\nrcore}{Core Nebulas Rank}
\newcommand{\nrext}{Extended Nebulas Ranks}
\newcommand{\dom}{{\; \texttt{dom}\;}}
\newcommand{\xpcomment}[1]{{\color{blue} #1}}

\setCJKsansfont[BoldFont = STHeitiSC-Medium]{STHeitiSC-Light}


\newtheorem{property}{特征}
%\addbibresource{reference.bib}

\begin{document}
\pagestyle{empty}
\renewcommand{\contentsname}{目录}
\renewcommand{\abstractname}{摘要}
\renewcommand{\refname}{参考文献}
%\renewcommand{\nomname}{术语表（按首字母排序）}
\renewcommand{\figurename}{图}
\renewcommand{\tablename}{表}
\renewcommand{\baselinestretch}{1.5}
\renewcommand{\appendixname}{附录}
\renewcommand{\proofname}{证明}

%\pagecolor{\pcolor}

\begin{titlepage}
  \begin{center}
    \vspace*{5.5cm}
    \includegraphics[scale=0.4]{../common/Nebulas.png}
    \vspace{0.5cm}


    \textbf{\huge{白皮书：NAX 机制设计}}

    \vspace{0.5cm}
    王卓尔
    \vfill
    2019年7月 \\
    版本号:0.0.1
    \textbf{}
  \end{center}

\end{titlepage}
\setcounter{page}{0}
%\thispagestyle{empty}
\tableofcontents
\newpage
\setcounter{page}{1}
\pagestyle{fancy}	
\vspace*{0.01cm}
\section{设计目标}

作为社区治理币，需要明确几个设计目标，最终NAX设计出来需要达到哪些目标。在确定一个系列目标的基础上，再进行机制的设计，所有机制的设计也将需要符合这个总的目标。为了使得机制设计简洁、清晰，一切与目标不相关的设计，都不应该随意纳入机制当中，以妨止最终偏离初衷。以下是几点关于治理币几个总的目标：

\begin{enumerate}[a.]
        \item 用于星云社区治理
        \item 星云社区贡献有效凭证
        \item NAX需要是符合通缩模型
        \item NAX有升值空间
        \item 销毁与增发达到一定程度的平衡
\end{enumerate}


\section{设计原则}

根据费雪公式：
\begin{equation}
M * V = T / P
\end{equation}
\(M\)是Token数量，\(V\)是Token流通速度，\(P\)是Token价格，而\(T\)是系统内总交易额。很好理解，等式两边其实算的都是以Token数量为计量的GDP。左边\(M * V\)是个数乘以流通速度等于GDP（Token计量），右边总GDP（法币计量）除以Token价格（法币计量）也等于GDP（Token计量）。

通过这个公式，我们不难推出(以后补上推理步骤）通过增加持币价值是最终有效的提升币价和使用价值的方式。

\subsection{减少NAS流通量}
\subsubsection{增加质押}
\subsubsection{增加地址数}
\subsubsection{减少交易所存量}
\subsection{增加持有NAS和NAX的动力（减少交易动力）}
\subsection{增加NAX使用和消耗场景 （供需平衡）}

% !TEX root = main.tex
\section{核心系统设计}

\subsection{质押NAS返NAX}
标准：
\begin{enumerate}
\item 四个周期（发行量）调整一次
\item 质押率 == 算力
\item 算力会影响难度
\item 衰减速率
\end{enumerate}

示例：
\begin{enumerate}[a.]

\item 一年减半（2,500,000)， 周期：100,000 (两周多）衰减一次\\
衰减系数：0.973\\
math.pow(0.973, 25） = 50\%

\item 一年减半（2,500,000)， 周期：50,000 (一周多）衰减一次\\
衰减系数：0.986\\
math.pow(0.986, 50） = 49.4\%

\item 两年减半（5,000,000)， 周期：100,000 (两周多）衰减一次\\
衰减系数：0.986\\
math.pow(0.986, 50） = 49.4\%

\item 两年减半（5,000,000)， 周期：50,000 (一周多）衰减一次\\
衰减系数：0.993\\
math.pow(0.993, 100） = 49.5\%

\end{enumerate}

新提议（质押消耗NAS) -- 感觉不是特别友好

如果我们需要维护一个质押率，达到一定的博弈平衡，可以添加，根据质押数量，消耗NAS的质押场景，也就是说，当质押分配到的NAX不合算的时候，用户可能会取消质押。当质押数变小的话，又会使得质押所得的NAX变得更多，所以又有人开始质押。

质押消耗的NAS可以被看作是BTC挖矿中消耗的电费。

收集的消耗的NAS会收集成为社区建设基金：Go Nebulas

公式如下：其中

        第\(i\)期用户\(j\)获得的NAX: \(K_{i,j}\)

        第\(i\)期用户\(j\)的质押量: \(P_{i,j}\)

        第\(i\)期用户\(j\)的质押时间: \(T_{i,j}\)

        第\(i\)期初始总增发量：\(C_i\)

        第\(i\)期增发比例： \(\lambda_i\)

        衰减系数：\(B\)

        h是质押高度总和，v是质押数量

\begin{equation}
  K_{i,j} = \frac{P_{i,j} T_{i,j}}{\sum_j P_{i,j} T_{i,j}} \lambda_i C_i
\end{equation}
\begin{equation}
  \lambda_i = f(\sum_j P_{i,j} T_{i,j})
\end{equation}
\begin{equation}
  C_i = C_0 B^i
\end{equation}

\subsection{质押返率\(\lambda\)}
\begin{enumerate}[a.]
  \item 时间。难度问题。早期质押， \(\lambda\)更高，正向量
  \item 预期？质押高度为质数， \(\lambda\)更高 x
  \item 运营活动。。。返回会越来越多？上一高度交易量越高， 此高度 \(\lambda\)更高 x
  \item 质押率（算力）越高， \(\lambda\)越高，质押率下降，难度下降，收益上升。质押率上升，难度上升。
\end{enumerate}

\begin{equation}
\lambda = (f(a, c, d) + g(h) b) / h
\end{equation}

\subsection{增发周期设定}
可以每100000高度分发一次，也就是高度满足以下性质：

H = h mod 100000 == 0

\subsection{取回质押策略}
满足性质：v越高，取回周期t越长，滞后返还, B是基本量级

B(v) =  floor(sqrt(v) - B) * t

\subsection{系统手续费}
每次增发的时候，新增发所得的4\% 转入NAX的专属项目基金。

矿工费？

固定手续费。发行费。铸币税。

项目团队的预期收益。


\section{使用场景设计}

各使用场景可以自行设定，奖励，销毁策略。

\subsection{Go Nebulas Platform: 生态激励，贡献，奖励（项目等）}

基金会资助Go Nebulas平台的300w NAS, 以及今后系统增发的部分NAS作为资助Go Nebulas平台经费，这些资金也将同时作为质押以获得NAX，作为社区贡献的附加奖励。

\subsubsection{奖励场景}

星云社区化项目制的开放性平台，社区贡献者可以由分为几个角色：社区开发者，社区推广者。社区贡献者通过Go Nebulas上做项目，将会获得NAS的工资奖励， 同时将会获得额外的NAX的社区贡献奖励/凭证，根据项目的优先级，定义不同比例的NAX返比。

\subsubsection{社区贡献者分类}

\begin{enumerate}[a.]
  \item 项目开发
  \item 运营/PR
  \item 市场拓展
  \item 拉新
  \item 基础设施：
  \item 主网
\begin{enumerate}[i.]
  \item NAS nano Pro
  \item 硬件钱包
  \item 新钱包对接
  \item 跨链合作
\end{enumerate}
\end{enumerate}

为了鼓励社区贡献者的积极性，我们初步确定以下的返NAX比例参数：

普通项目 ：S * x * N (NAS reward)

优先项目（基础设施）：T * x * N (NAS reward)

(s < t)

\subsubsection{消耗场景}
\paragraph{创建提案（消耗相应的NAX）}
\paragraph{开发提案（开发时会销毁等比的NAX, 提前完成的工作，节省销毁的数量）}
\paragraph{投票通过提案和结果，需要销毁后NAX（通过或以获得120\%返还，失败可以获得110\%的返还）}

\subsubsection{推进步骤}
\paragraph{补发过去参与过贡献的人的奖励？}
\paragraph{以后获得项目资金的人将获得相应NAX奖励}
\paragraph{增加GN邀请奖励，受邀请的贡献者获得的NAX后，邀请人会获得额外10\%的NAX}

\subsection{投票（Voting）提升社区参与度}
为了鼓励社区使用NAX作为社区治理的工具，以及鼓励大家社区治理的参与度，我们将设置一些投票场景，部分场景投票将做会有返还奖励。
投票的场景

\subsubsection{节点竞选}
\begin{enumerate}[a.]
  \item 从节点候选人里可以选出节点
  \item 投票销毁后会返还（若当选，获得120\%返还，失败有110\%返还）
\end{enumerate}

\subsubsection{理事竞选}
\begin{enumerate}[a.]
  \item 从候选人里选出相应的理事
  \item 投票销毁后会返还（若当选，获得120\%返还，失败有110\%返还）
\end{enumerate}

\subsubsection{销毁投票}
\begin{enumerate}[a.]
  \item 利益相关的投票，投票将没有返还
  \item 如投票通过销毁3500w部分token，这个投票，将不会返还
\end{enumerate}

\subsection{当选节点 ---- 一次性销毁门槛}

\subsubsection{参与节点竞选}
需要销毁1000W NAX，具体参与需要再调整

\subsubsection{成为节点}
被选中成为出块节点后，参与到节点出块，需要质押NAS(无NAX返还）


\subsection{NAS nano Pro \& Explorer \& DEX等各类生态平台的上币费/手续费}
随着Next DAO的推进以及社区治理的前进，社区里将会出现越来越多的Token和治理尝试。这些币种都将需要相应的工具支持，所需要上NAS nano Pro和Explorer的需求。资源空间有限的情况下，我们可以采取增加上币费的需求，比如需要在NAS nano Pro和Explorer上币的NRC20需要缴纳 500w NAX（参数可调）的上币费。其中20\%归集给NAS nano Pro和Explorer管理团队 （开发者和运营者）80\%会被销毁

\subsection{社区预留NAS销毁计划（单独的方案）x}
与其将3500w NAS一次投票销毁，其实可以把销毁做成一项长期的社区投票活动，由社区来决定这个事情的发生。

销毁细节：

每个自然月1号发起一次投票销毁社区预留剩余NAS总量的 \(\alpha\) \% ， \(\alpha\) \% 是当前NAS 质押率占流通量的份额。

投票销毁通过细节：

\begin{enumerate}[a.]
  \item 投票需要满足 \(\alpha\) \%的现行NAX总量，才合格
  \item 支持销毁的比例超出50\%
  \item 投票销毁，只返还50\%
\end{enumerate}

nax的发行在nextDAO和nebulas.io同时可以看到

Nrc20

DAI, bridge / 其他公链的币可以在nextDAO平台上做映射


技术下沉主网需要规划

\section{NAX分析}
由于每期增发池中未发放的部分滚入下一期增发池中，因此总发行量为固定值
\begin{equation}
  \sum_{i,j} K_{i,j} = \sum_i C_0 \mu^i = \frac{C_0}{1-\mu}
  \end{equation}
  令此上界为100亿（\(10^{10}\)），可解出\(C_0 = 10^{10}(1-\mu) = 1.9\times10^7\)。

\subsection{发行量分析}


\subsection{反作弊分析}

%\printbibliography
\newpage
\bibliography{reference}

\newpage 
\begin{appendices}
\section{Change Log}
\begin{itemize}
\item{0.0.1} Release. 2019.8
\item{0.0.2} Release. 2019.11
\end{itemize}

\end{appendices}

\end{document}
